% ----------------------------------------------------------------------------
%                                 CONCLUSÃO
% ----------------------------------------------------------------------------

\chapter{CONSIDERAÇÕES FINAIS}
\label{chap:conclusao}

O desenvolvimento de tecnologias aplicadas à área da saúde tem se tornado cada vez mais relevante no suporte ao acompanhamento clínico e na tomada de decisões médicas. Nesse contexto, este trabalho apresentou a proposta de adaptação do sistema Painel Parkinson para o ambiente \textit{desktop}, buscando ampliar a acessibilidade, melhorar a experiência de utilização e reduzir limitações associadas ao acesso exclusivamente via navegador \textit{web}.

Ao longo do trabalho, foram apresentados os conceitos relacionados à Doença de Parkinson, aos sistemas de apoio à decisão clínica e à utilização de \textit{dashboards} para visualização de dados biomédicos. Além disso, foram discutidos os aspectos arquiteturais e tecnológicos envolvidos na construção da solução, incluindo a utilização das tecnologias Electron.js, VueJS, Flask e PostgreSQL.

A proposta da aplicação \textit{desktop} visa proporcionar maior praticidade para os profissionais da saúde durante o acompanhamento clínico dos pacientes, permitindo acesso rápido às informações, redução da dependência de navegadores e potencial melhoria no desempenho da aplicação em ambiente local.

Também foram apresentadas as principais interfaces do sistema, demonstrando a organização dos módulos clínicos, \textit{dashboards}, a estruturação da escala MDS-UPDRS e os mecanismos de visualização de dados relacionados à Doença de Parkinson. Os testes preliminares indicam que a abordagem \textit{desktop} possui potencial para contribuir positivamente na usabilidade e na eficiência do sistema no contexto clínico.

Como continuidade do projeto, pretende-se realizar melhorias relacionadas à integração com o banco de dados, otimização de desempenho, segurança da informação e ampliação dos testes de usabilidade com profissionais da saúde. Adicionalmente, a reintrodução do processamento local de sensores inerciais (acelerometria e giroscopia) e das escalas psicológicas e emocionais complementares será abordada em trabalhos futuros. Além disso, futuras versões poderão incorporar novos recursos voltados à análise de dados clínicos e monitoramento contínuo dos pacientes.

Dessa forma, espera-se que o presente trabalho contribua para a evolução do Painel Parkinson e para o fortalecimento de soluções tecnológicas aplicadas ao suporte à decisão clínica no contexto da Doença de Parkinson.

Por fim, em conformidade com as diretrizes de transparência acadêmica, ressalta-se que ferramentas de Inteligência Artificial, como o Google Gemini, foram utilizadas de forma estritamente assistiva ao longo da redação deste documento, auxiliando na revisão gramatical e na formatação tipográfica. Todas as modelagens, códigos e análises clínicas apresentadas permanecem de autoria intelectual exclusiva dos pesquisadores.